\documentclass[sigconf,nonacm]{acmart}
\setcopyright{none}
\settopmatter{printacmref=false,printfolios=true}
\usepackage{tabularx}
\raggedbottom
\hypersetup{keeppdfinfo,pdfauthor={Angelo Dalli},pdftitle={Separating Grounding, Execution and Attribution in Foundation-Model-Assisted Banking Audits}}
\begin{document}
\title[Separating Grounding, Execution and Attribution]{Separating Grounding, Execution and Attribution in Foundation-Model-Assisted Banking Audits}
\author{Angelo Dalli}
\email{angelo.dalli@umnai.com}
\affiliation{\institution{UMNAI}\city{}\country{}}
\renewcommand{\shortauthors}{Angelo Dalli}
\begin{abstract}
Auditing a foundation-model-assisted banking decision means keeping three questions apart: whether the cited source supports the value that was used, whether the policy ran correctly on that value, and whether the final action came from the policy alone or involved a later human intervention. This paper defines an executable audit contract with seven assertions. Each has its own evidence obligation and is judged supported, contradicted or unresolved. A reproducible synthetic merchant-onboarding suite shows that these judgments can differ: a rule can replay correctly on a misread commitment, and a recommendation can be stable while the source reading remains unsettled.

In archived outputs from a model run without computational tools, the same grounding-prerequisite failure appears in two fixtures. The model supports the recorded commitment, but local hash recomputation shows that the citation digest does not match the retained contract. Commitment grounding should therefore be unresolved. In one response, the model marks lineage as contradicted and commitment grounding as supported. Keeping these labels separate makes the bypass of citation validation visible. The frozen suite tests conformance to the authored contract, and the archived outputs can be checked offline. The contribution is this separation of evidence obligations and the concrete failure it exposes. Tests on independent documents, and any benefit to human reviewers, remain future work.
\end{abstract}
\ccsdesc[500]{Computing methodologies~Artificial intelligence}
\ccsdesc[300]{Applied computing~Banking}
\ccsdesc[300]{Information systems~Data provenance}
\keywords{claim verification, foundation models, decision provenance, evidence grounding, auditability}
\maketitle
\hypersetup{pdfauthor={Angelo Dalli}}
\section{Introduction}

A banking decision may combine document interpretation, transaction aggregates, policy logic and human review. Each stage creates a different audit question. A reproducible rule trace shows how recorded inputs produce a recommendation, yet it does not establish that those inputs faithfully represent a cited contract. A recommendation and a final action may coincide even when a human intervened. Treating these questions as a single judgment obscures the evidence required for each.

In a hypothetical merchant-onboarding workflow, a foundation model extracts cancellation terms, a service supplies dispute counts, a policy recommends a route and a reviewer may intervene. Reviewability requires records across this process~\cite{cobbe2021reviewable}.

A foundation model can extract the inputs being audited or act as the auditor; these are distinct roles. This paper examines both on synthetic fixtures. It contributes an executable contract separating seven assertions, a suite demonstrating divergent outcomes, and archived outputs supporting failure analysis. The contract preserves supported execution when grounding is unresolved and distinguishes a stable recommendation from a settled interpretation. Two tool-free responses endorse an input despite an invalid citation; separate lineage and grounding labels expose the prerequisite failure.

\paragraph{Relation to prior work.}
Faithfulness concerns whether explanations reflect computation~\cite{jacovi2020}; model rationales may omit influencing factors~\cite{turpin2023}. Lyu et al.~separate language-model translation from symbolic execution~\cite{lyu2023}. Decision provenance~\cite{singh2019} and PROV~\cite{prov2013} record derivations. The present contract checks specific assertions within a retained decision record.

FEVER distinguishes support, refutation and insufficient evidence~\cite{thorne2018fever}; AIS evaluates attribution to identified sources~\cite{rashkin2023measuring}. Assurance cases connect claims, evidence and assumptions; SACM represents unsupported and defeated claims~\cite{sacm2023}. Here these principles distinguish grounding, replay, intervention history and sensitivity over declared interpretations. Reporting resolution alongside errors follows selective prediction~\cite{geifman2019}; abstention reflects unavailable evidence or checking capability.

\section{Executable claim contract}

\subsection{Records, assertions and outcomes}

Let \(z\) denote consumed features and \(s\) the policy and remaining consumed state. The policy recommends \(y=d(z,s)\); the final action \(a\) may additionally depend on an intervention. Each claim kind \(k\) specifies evidence obligations \(D_k\), assumptions \(A_k\) and a check \(K_k\). A supplied claim has an identifier, kind and asserted value \(v\). For the retained record \(E\),

\begin{equation}
K_k(E,v;A_k)=(\ell,r,B),\qquad
\ell\in\{\mathrm{S},\mathrm{C},\mathrm{U}\},
\end{equation}

where \(r\) is a reason and \(B\) lists evidence references. \emph{Supported} (S) means the assertion meets its obligations; \emph{contradicted} (C) requires decisive counter-evidence; \emph{unresolved} (U) covers missing or conflicting prerequisites, unsupported language or disabled checks. Records retain sources, citations, consumed inputs, policy, trace, actions, intervention history, completeness declarations, candidates and claims. Table~\ref{tab:claims} specifies the seven implemented scopes; they are selected obligations for this workflow, not an exhaustive banking-audit taxonomy.

\begin{table}[t]
\caption{Seven audit assertions and their evidence obligations. Execution, final-action attribution and dispute sufficiency concern consumed inputs; they do not inherit source grounding.}
\label{tab:claims}
\small
\begin{tabularx}{\columnwidth}{@{}lX@{}}
\toprule
Assertion & Evidence required by the contract \\
\midrule
Contract lineage & Cited identity, version and byte digest match the retained contract. \\
Commitment grounding & Valid cited contract supports the asserted penalty-free exit time. \\
Dispute grounding & Valid transaction source supports an equal asserted fraction in the specified window. \\
Conditional execution & Both recorded predicates and the recommendation match exact replay. \\
Final-action attribution & Valid execution, matching action, no intervention and stipulated complete history. \\
Candidate stability & Recommendations are invariant over the declared candidate set. \\
Sufficient dispute reason & Recorded disputed and eligible counts satisfy $50u\geq n$. \\
\bottomrule
\end{tabularx}
\end{table}

For a true lineage assertion, matching retained bytes and versions yield S; an observed mismatch yields C; an absent or ambiguous artifact yields U. \emph{Grounding} means support from the cited, integrity-validated source, rather than agreement with an uncited passage. A digest mismatch therefore contradicts lineage but leaves grounding unresolved: the cited source is unusable for that assertion. With valid citations, a unique interpreted value supports an equal assertion and contradicts a different one. Missing, unsupported or equally authoritative conflicting interpretations yield U. Dispute grounding compares fractions by exact cross multiplication.

Conditional execution checks both predicates and the recommendation against replay on consumed inputs. An explicit intervention, including a same-valued one, contradicts policy-only final-action attribution. A final-action mismatch with valid replay also contradicts that assertion. Without decisive counter-evidence, an invalid trace or incomplete history yields U. Supported attribution requires valid execution, a matching action, no intervention and stipulated complete, trustworthy history. Unresolved source grounding therefore leaves replay independently checkable, while failed replay prevents supported attribution.

Lineage, stability and dispute-sufficiency claims have Boolean asserted values: a false-valued assertion is supported when its predicate is false. Claim text must express the structured proposition consistently. Hashes establish consistency with retained bytes, not source authenticity; grounding checks fidelity to retained text, not real-world truth. Historical attribution additionally depends on trustworthy records.

\subsection{Grounding and decision sensitivity}

The same consumed commitment and trace can accompany different source interpretations. Replay cannot distinguish their grounding conditions; conflicting interpretations can still yield the same recommendation.

For a declared nonempty candidate set \(Z\), hold remaining inputs and policy fixed and enumerate \(\mathcal{Y}(Z;s)=\{d(z,s):z\in Z\}\). A singleton outcome set supports a true stability assertion; multiple outcomes contradict it. Empty or unavailable candidates yield U. Stability establishes invariance only over \(Z\); it does not resolve the interpretation or establish candidate completeness. This implementation varies commitment alone. Candidate generation, joint uncertainty and real-world causal effects~\cite{karimi2021} are outside that calculation.

\section{Synthetic regression study}

\subsection{Fixtures and authored references}

The frozen test suite release for this paper consists of 100 episodes and 700 claims, plus 8 diagnostic episodes and 56 claims. A total of 20 deterministic parameter combinations (families) each have 5 variants: clean evidence, extraction error, source gap, ambiguity and intervention/history issue. Additionally, 5 transaction-count pairs, 6 commitment durations, 3 framework durations and 3 numeric contract templates cycle together. One transaction template fixes the window; count errors add 5 consumed pairs. All merchants and documents are fictional.

The illustrative onboarding screen uses disputed count \(u\), eligible count \(n\) and penalty-free commitment \(c\) in months:

\begin{equation}
d(u,n,c)=
\begin{cases}
\text{manual review}, & 50u\geq n\ \lor\ c<6,\\
\text{standard processing}, & \text{otherwise}.
\end{cases}
\end{equation}

Here \(0\leq u\leq n\), \(n>0\) and \(c\geq0\) are integers. A dispute rate of at least \(2\%\) triggers review. The commitment threshold is an invented screen, not an empirically validated risk indicator. Framework duration and penalty-free exit time are distinct fields.

Fourteen non-clean subtypes inject framework-as-commitment substitutions, disputed-count errors, missing artifacts, conflicting contracts, digest/version faults, spelled-out numbers, overrides, incomplete histories and incorrect traces. Four of twenty extraction errors change the recommendation. The original suite excludes combined faults and richer document language. Its 73 singleton candidate sets, four empty sets and 23 pairs cover declared-candidate stability. All pairs straddle the six-month boundary; the additional diagnostics include same-side and three-candidate sets.

Reference labels and rationales derive from the authored constructions. The generator does not import the checker, but both encode the same semantics. Eight families form the development subset (40 episodes); twelve form the reporting subset (60 episodes, 420 claims). Parameter values and templates recur across subsets, and reporting cases were inspected during development. Consequently, agreement measures conformance to the authored rules, not blind generalisation, with the dataset card recording the construction~\cite{gebru2021}.

\subsection{Checker configurations and reporting}

All checker configurations receive identical sources, citations and records. \emph{Execution} checks replay, final-action attribution, candidate stability and dispute sufficiency. \emph{+Identity} adds contract identity, version and SHA-256 verification. \emph{+Grounding} compares asserted values with sources interpreted by full-match parsers for three numeric contract templates and one transaction template. Disabled checks and unsupported wording yield U. Exact integer arithmetic avoids threshold drift.

Reported measures are three-class reference agreement, resolution (S or C), and \emph{false support}: predicted S where the authored reference is C or U. False support is reported per supported prediction (endorsement). It concerns unjustified endorsement under the contract; it does not establish that the underlying text or merchant decision is false. Reference U denotes insufficient evidence, while checker U can also reflect a capability limit.

The artifact reports confusion matrices and scope, variant and family breakdowns. Identifiers are used only to join outputs; identifiers, split, variant and reference labels are excluded from decision features. Counts are descriptive because fixtures share structure. Scripts regenerate fixtures and rescore outputs offline. Contract, reporting labels and archives are hashed while diagnostics are held separate.

\subsection{Archived model procedures}

Three foundation models were used in the evaluation: GPT-6.1 Sol Ultra, Grok 4.7 and Claude Fable 5.1. Sol Ultra was evaluated using the artifact-authoring assistant system. Four declared fresh 17-episode contexts cover 60 reporting episodes and eight diagnostics, without reference labels, rationales or scenario tags. The worker archive declares no inherited conversation history and prohibits memory or project lookup. Parent configuration was inherited; isolation from persistent memory, project context or instruction state was not independently verified. Three batches mix reporting and diagnostics; 37/60 reporting episodes share a context with a same-family variant. A frontend supplies citation presence, kind, version and digest comparisons. Integrity is precomputed; prerequisite application, source interpretation and policy reasoning remain model tasks. Paired rendered-evidence audits of twelve episodes use fresh six-episode batches: PNG cards replace passages while populated records and integrity observations remain visible. A rendering script in the supplement reproduces all 25 PNGs byte for byte with Pillow 12.3.0.

Sol Ultra uses one attempt per batch without repair. Invalid/\allowbreak{}missing batches receive U placeholders counted incorrect even against reference U; denominators remain fixed. Instructions, literal responses and self-reported activities are archived.

Grok predictions cover 100 original episodes and eight diagnostics; reporting and diagnostic outputs are rescored separately against authored references. The supplied account describes per-episode judgment with hashing and exact-replay tools. Native prompts and execution traces are absent, so the inference procedure cannot be examined.

Fable reports fresh instances per episode under two conditions: no computational tools, or Bash restricted to SHA-256 and integer arithmetic. Each audits 60 reporting episodes and eight diagnostics. Prompts, 136 literal responses and conversion files support rescoring; no response failures are recorded. Fable additionally extracts 55 fields from 25 unique documents (20 contracts, five transaction summaries). Text and rendered-image conditions each use 25 declared fresh calls without populated audit records. Images have mild synthetic rotation, blur and speckle. Prompts, condition wrappers, images, outputs and reference fields are archived.

\section{Results and failure analysis}

\subsection{Grounding endorsements with invalid citations}

The same grounding-prerequisite failure occurs in two Fable responses without computational tools (Table~\ref{tab:failures}). Both support commitment grounding despite invalid all-zero citation digests; local SHA-256 recomputation confirms both mismatches. In response A, the model contradicts lineage yet supports grounding, treating the digest defect as a separate issue. In response B, it abstains on lineage yet supports grounding. Both bypass the grounding prerequisite: a readable value in the retained text does not validate its citation.

\begin{table}[t]
\caption{One grounding-prerequisite failure observed twice in the tool-free Fable run. Both citation digests mismatch; required grounding is U. S/C/U denote supported/contradicted/unresolved. Episode identifiers and responses are retained in the artifact.}
\label{tab:failures}
\small
\begin{tabularx}{\columnwidth}{@{}Xccc@{}}
\toprule
Response & Lineage & Grounding & Required\\
 & predicted & predicted & grounding\\
\midrule
A & C & S & U\\
B & U & S & U\\
\bottomrule
\end{tabularx}
\end{table}

The hash/arithmetic run's zero false support does not establish that tools prevented this error. Hashing and arithmetic changed together, and separate calls had unrecorded sampling. The tool-free run agrees with 359/420 references, resolves 331 claims and has 2/277 false support. Its other 59 disagreements are abstentions: 55 lineage and four dispute grounding. Abstaining without hashing avoids an unsupported integrity claim. For dispute grounding, four responses cite unavailable hashing when abstaining, 54 give S/C judgments and two lack the source. Diagnostic agreement is 43/56, with 0/23 false support. Each fixture has one attempt per condition; no plausible wrong digest or combined fault was tested. Only Fable had a tool-free condition; equivalent runs for the other systems are absent. These observations establish one failure mode, without comparing model susceptibility or estimating its frequency on realistic citations.

\subsection{Distinct outcomes across audit scopes}

\paragraph{Extraction and execution.}
All twelve reporting extraction errors contradict grounding while preserving execution; two change the recommendation. For example, a contract specifies a three-month penalty-free exit and a 24-month framework. Consuming \(c=24\) at \(u=25,n=1500\) yields standard processing. Lineage is supported, commitment grounding contradicted and replay supported. A matching intervention-free final action with complete history supports attribution to that replay. Correcting \(c\) to three triggers review. Supported computation therefore leaves the extraction error intact.

\paragraph{Grounding and stability.}
Twelve ambiguity episodes leave commitment grounding unresolved: seven are stable over \(\{3,9\}\) and five are sensitive. At \(25/1500\), those alternatives produce different routes; at \(25/1200\), disputes trigger review under either reading. Two further source-gap conflicts complete fourteen conflicting-contract cases. Stability can thus coexist with unresolved grounding.

\paragraph{Stability and dispute sufficiency.}
Table~\ref{tab:diagnostics} extends the original straddling pairs. Each set has positive and negative stability assertions, giving eight diagnostic episodes. At \(u=25,n=1500\), commitment grounding is U and true dispute-sufficiency assertions are C. The checker agrees with all 56 diagnostic labels. For this screen, with \(P=(50u\geq n)\), stability holds when \(P\), all candidates are below six, or all are at least six. It reduces to \(P\) only for straddling sets. This policy-specific identity illustrates why candidate stability and dispute sufficiency have distinct evidence obligations.

\begin{table}[t]
\caption{Candidate diagnostics at \(u=25,n=1500\). K(T)/K(F) evaluate true/false stability assertions. Candidate sets are declared, not claimed exhaustive.}
\label{tab:diagnostics}
\small
\begin{tabularx}{\columnwidth}{@{}Xccc@{}}
\toprule
Candidate commitments & Outputs & K(T) & K(F)\\
\midrule
$\{6,9\}$ & standard & S & C\\
$\{3,5\}$ & review & S & C\\
$\{3,6,9\}$ & both & C & S\\
$\{6,9,12\}$ & standard & S & C\\
\bottomrule
\end{tabularx}
\end{table}

\paragraph{Action and attribution.}
Three same-valued interventions contradict policy-only final-action attribution despite matching recommendation and action. Three faulty-trace episodes instead yield contradicted execution and unresolved attribution. Equal action values establish neither intervention-free history nor valid execution.

\subsection{Conformance counts and paired modalities}

The full checker agrees with 418/420 labels, resolving 386 (331 S, 55 C; Table~\ref{tab:results}). Its 34 U include two supported-reference commitments worded as \emph{eighteen}, outside the digit-only grammar. The misses reflect the declared parser boundary. Zero false support applies to these fixtures.

\begin{table}[t]
\caption{Checker configurations on 420 authored claims. Disabled checks abstain. Generator and checker share semantics; these are conformance counts. FS/S is false support per endorsement.}
\label{tab:results}
\small
\begin{tabularx}{\columnwidth}{@{}Xrrr@{}}
\toprule
Enabled checks & Agreement & Resolved & FS/S\\
\midrule
Execution & 264/420 & 232/420 & 0/193\\
+ Identity & 322/420 & 290/420 & 0/247\\
+ Grounding & 418/420 & 386/420 & 0/331\\
\bottomrule
\end{tabularx}
\end{table}

Archived Sol Ultra, Grok predictions and Fable with hashing and arithmetic each agree with 420/420 authored labels: 333 S, 55 C and 32 U, with 388 resolved and 0/333 false support. All reach the authored resolution ceiling of 388/420, so this subset cannot distinguish them by agreement. Unmatched procedures and unverified provenance further prevent comparative performance claims. The 32 U are required abstentions: 20 commitment grounding, six attribution and two each lineage, dispute grounding and stability; none denotes a failed response. Each also matches 56/56 diagnostics (36 S, 12 C, eight U), with 0/36 false support. A supplemental inventory records prompts, responses and procedures; Grok predictions permit rescoring only.

Sol Ultra's paired rendered-evidence audits each match 84/84 labels (24/24 grounding; 68 S, seven C, nine U; 0/68 false support), without batch failures. In the separate extractor role, Fable's source-only extraction matches 55/55 fields and 25/25 documents per modality, without framework-as-commitment errors. Literal quotations match source substrings in 20/25 text and 21/25 image responses; whitespace normalisation yields 25/25 and 24/25. One image quotation replaces line breaks with a separator. Field correctness does not ensure verbatim attribution. These are paired views of 25 templated documents, not 50 independent documents or realistic scans.

\section{Discussion and limitations}

Supported replay can coexist with contradicted grounding, and stability with unresolved interpretations. Contradiction could prompt claim correction; unresolved status could prompt evidence collection or qualification. These are intended uses, without measured reviewer benefit, and prescribe no merchant action. References and checker share assumptions; reporting fixtures were inspected during development. Agreement establishes internal consistency. Minimal templates, correlated parameters, selected faults and stipulated histories limit generalisation. Parsing is bounded, hashes do not authenticate documents, and candidates may omit valid interpretations.

The model results are archived observations rather than a controlled comparison. Shared provenance across Sol Ultra's artifact construction and auditing adds dependence beyond generator/\allowbreak{}checker conformance; sibling variants also co-occur within batches. The Fable operator had seen references. Fresh-context isolation and tool restrictions are reported by the foundation models themselves rather than independently verified. Fairness, compliance and causal actionability were not evaluated.

\paragraph{Future evaluation.}
Independent challenge testing remains future work: new documents, plausible wrong digests and combined faults; blinded annotation and adjudication independent of checker authors; labels frozen before inference under matched prompts and tools. Annotator disagreement should accompany model results. A subsequent reviewer study would test fault-identification benefit.

\section{Conclusion}

Grounding, policy replay and final-action attribution require separate evidence. On authored synthetic fixtures, the contract exposes one grounding-prerequisite failure in two responses, including an inconsistent lineage/grounding judgment. Released records, labels, outputs and the scorer permit offline examination. The results establish scope separation and fixture conformance; independent audit evaluation remains future work.

\paragraph{Data and code availability.}
The synthetic dataset, checker, authored references and archived outputs accompany this manuscript as ancillary files. Appendices~\ref{sec:supp-contract}--\ref{sec:supp-reproducibility} give the supplementary methods, run records and reproduction details.

\clearpage
\balance
\bibliographystyle{ACM-Reference-Format}
\bibliography{references}
\clearpage
\nobalance
\appendix
\section{Supplementary contract specification}
\label{app:contract}
\label{sec:supp-contract}

The released contract makes a claim about a retained decision record, rather than about an unobserved banking process. Its observable episode contains source texts; citations with source identity, version and SHA-256; consumed inputs; policy configuration; both recorded trace predicates; policy recommendation; final action; an optional intervention; an action-history completeness declaration; candidate commitments; and seven structured assertions. Each assertion contains a claim identifier, kind, asserted value and human-readable rendering. The checker evaluates the structured proposition; arbitrary prose entailment is outside its scope. Each output retains a label, check description and evidence references.

The source texts are retained UTF-8 strings. Citation digests cover the exact encoded string, including its line breaks. The checker resolves each cited identity against the retained list. A present source with the wrong kind, version or digest is decisive counter-evidence for a true lineage assertion. A missing, ambiguous or malformed reference is unresolved unless another reference supplies a decisive mismatch. A matching digest establishes consistency with retained bytes. It does not authenticate a document or establish its completeness, authority or truth.

The following procedure summarises the grounding prerequisite. \texttt{unique\_value} includes the supported-language check and the absence of equally authoritative conflicting readings. Commitment grounding uses an integer value; dispute grounding uses exact equality of fractions in the declared observation window.

\begin{quote}
\small
\begin{verbatim}
grounding(assertion, citations, sources):
  status, validated = check_citations(
      citations, sources)
  if status != supported:
      return unresolved
  value = unique_value(validated)
  if value is unavailable:
      return unresolved
  if equal(value, assertion):
      return supported
  return contradicted
\end{verbatim}
\end{quote}

\paragraph{Contract lineage.}
The assertion concerns every cited contract's identity, version and bytes. For a true assertion, all required matches yield S, a decisive mismatch yields C, and an unavailable comparison yields U. A false-valued lineage assertion reverses S and C when the predicate can be decided; U remains U.

\paragraph{Commitment grounding.}
All required citations must validate before their text can support the claimed penalty-free exit time. The bounded checker full-matches one of three numeric contract grammars. An unsupported grammar, unusable citation or conflicting valid values yields U. A unique value equal to the asserted commitment yields S; a unique different value yields C. Framework duration is a separate quantity. The claim does not ask whether using a different value would alter the recommendation.

\paragraph{Dispute grounding.}
Validated transaction sources must supply valid nonnegative disputed counts and positive eligible counts, with disputed count no greater than eligible count, in the window \texttt{2026-01-01/2026-01-31}. For source counts $(u,n)$ and asserted counts $(u',n')$, equality is $un'=u'n$. Thus the assertion is about the fraction, not separate equality of the two counts. Unusable or conflicting source fractions yield U. A unique equal fraction yields S and a different fraction yields C.

\paragraph{Conditional execution.}
The supported policy is exactly \texttt{toy-v1}, with dispute threshold $1/50$ and commitment threshold six months. Required inputs are integers satisfying $0\leq u\leq n$, $n>0$ and $c\geq0$. Both trace predicates must be present as Booleans: $50u\geq n$ and $c<6$. The recorded recommendation and asserted recommendation must agree with replay. A valid but inconsistent trace or recommendation yields C; unavailable inputs, an unsupported policy or malformed required fields yield U. This check uses consumed inputs independently of their source grounding.

\paragraph{Policy-only final-action attribution.}
An assertion naming a different retained final action is C. An explicit intervention is C even when its before and after values coincide. A final action differing from a replayable policy output is also C. In the absence of decisive counter-evidence, missing completeness information or an unavailable or inconsistent execution record yields U. S requires a matching action, explicit absence of an intervention, valid execution and stipulated complete, trustworthy history. Completeness is a fixture assumption; the implementation does not establish completeness of real logs.

\paragraph{Candidate stability.}
For the declared nonempty set of nonnegative integer commitments, remaining consumed inputs and policy are held fixed. The checker enumerates the recommendations and compares the resulting singleton/non-singleton predicate with the asserted Boolean. Agreement yields S and disagreement C. Invalid or unavailable candidates, or incomplete replay, yield U. The result is conditional on the declared set and makes no claim that all reasonable interpretations have been enumerated.

\paragraph{Sufficient dispute reason.}
The checker compares the asserted Boolean with $50u\geq n$ under the supported policy and recorded dispute counts. This assertion does not require a commitment value, source grounding or intervention-free final action. Invalid required counts, policy or assertion yield U. Otherwise predicate agreement yields S and disagreement C.

These obligations permit combinations that a single audit verdict would conceal. A readable but invalidly cited contract can yield lineage C and commitment grounding U while execution remains S. A grounded extraction error yields grounding C and execution S. A same-valued intervention can yield execution S and policy-only attribution C. Conflicting source readings can yield grounding U and candidate stability S. These are statements about different propositions, not contradictory verdicts on one proposition.

\section{Fixture construction and authored references}
\label{app:fixtures}

Release 1.2.2 preserves the original main corpus at 100 episodes and 700 claims, with a separate eight-episode, 56-claim diagnostic extension. The generator uses deterministic cyclic assignments over twenty families rather than sampling banking records or taking a full Cartesian product. With zero-based family index $i$, transaction pairs cycle with $i\bmod5$, commitments with $i\bmod6$, and framework durations and numeric contract grammars with $i\bmod3$.

The five transaction pairs are $15/800$, $20/1000$, $25/1200$, $25/1500$ and $45/2000$. Commitments are 3, 5, 6, 9, 12 and 18 months; framework durations are 24, 36 and 48 months. The transaction window is fixed. Every family has clean, extraction-error, source-gap, ambiguity and intervention/history variants. The public development partition contains eight families and forty episodes; the reporting partition contains twelve families and sixty episodes. All variants of one family remain together. Values and language recur across the partitions, both of which were inspected during implementation.

Table~\ref{tab:supp-faults} records the non-clean constructions. Extraction faults enter the consumed inputs before policy execution. Except for the explicit trace-fault subtype, replay therefore faithfully executes the recorded values. Four of twenty extraction errors change the recommendation relative to their clean counterpart, including two of twelve in the reporting partition. Fault frequencies are design choices rather than estimated deployment frequencies.

\begin{table}[t]
\caption{Authored non-clean subtypes. Development and reporting counts describe the original corpus; the separate diagnostics are excluded.}
\label{tab:supp-faults}
\small
\begin{tabularx}{\columnwidth}{@{}Xrr@{}}
\toprule
Construction & Development & Reporting\\
\midrule
Framework substituted for commitment & 4 & 6\\
Disputed count increment & 4 & 6\\
Missing contract & 2 & 2\\
Missing transactions & 2 & 2\\
Conflicting contracts (source gap) & 1 & 2\\
Invalid contract digest & 1 & 2\\
Wrong contract version & 1 & 2\\
Unsupported worded contract & 1 & 2\\
Conflicting commitments, sensitive & 3 & 5\\
Conflicting commitments, stable & 5 & 7\\
Changed-output intervention & 2 & 3\\
Same-valued intervention & 2 & 3\\
Incomplete history & 2 & 3\\
Incorrect trace predicate & 2 & 3\\
\midrule
Total non-clean episodes & 32 & 48\\
\bottomrule
\end{tabularx}
\end{table}

The numeric contract grammars express penalty-free cancellation time separately from service-framework duration, with one of three fixed word orders. The transaction grammar supplies an observation-window line followed by eligible and disputed counts. The worded-contract subtype supplies the human-readable commitment \emph{eighteen} but falls outside the checker grammar. Its authored grounding label is S, while the bounded checker returns U. This accounts for both reporting disagreements of the full checker.

References and rationales derive from the constructions. The generator does not import the checker, but both encode the same contract semantics. The reference file therefore supplies an authored answer key, rather than independent annotation. Hidden fictional facts are not used to resolve an absent or conflicting source. The main reporting key has 333 S, 55 C and 32 U labels. The 32 required abstentions comprise twenty commitment-grounding claims, six final-action-attribution claims, and two claims each for lineage, dispute grounding and candidate stability. They are evidence-limited outcomes rather than response failures.

The original corpus declares 73 singleton candidate sets, four empty sets and 23 pairs. Its ambiguity pairs straddle the six-month threshold. The separate diagnostic extension fixes $u=25,n=1500$ and declares $\{6,9\}$, $\{3,5\}$, $\{3,6,9\}$ or $\{6,9,12\}$. Each set has one true-valued and one false-valued stability assertion, giving eight episodes. These paired constructions test asserted polarity and separate stability from dispute sufficiency. They do not add independent documents or sampled populations, and their counts are never pooled with the 420 reporting claims.

\section{Grounding-prerequisite failure records}
\label{app:failures}

Response A is archived as \texttt{case\_96101184a60febe1}; response B as \texttt{case\_91bf94e28c04fb8c}. Both retained contract texts state a nine-month penalty-free cancellation time and a 24-month framework duration. Their contract citations specify 64 zero digits as the SHA-256 digest. Computing SHA-256 over the exact retained UTF-8 bytes confirms a mismatch in both episodes. The appearance of an all-zero digest alone is not a mathematical proof of that mismatch.

In A, consumed counts are $45/2000$ and the recorded recommendation is manual review. The model emits lineage C and commitment grounding S. Its grounding reason states: ``The digest problem in c1 is a lineage defect assessed separately''. The record therefore exposes a dependency error after the model has itself recognised a citation defect: under the contract, the retained passage is not usable evidence for that grounding assertion.

In B, consumed counts are $15/800$ and the recorded recommendation is standard processing. The model emits lineage U and commitment grounding S. Its lineage reason says the digest ``cannot be checked without computation''. It nevertheless accepts the readable commitment as grounded. The two cases share one failure mode: commitment grounding is endorsed without satisfying its citation-validation prerequisite. The required grounding label is U in both; neither case shows a dispute about the nine-month wording.

The raw responses preserve every emitted claim, reason, evidence reference and activity self-report under the Claude archive's \texttt{raw/audit\_notools/} directory. These are two observations on conspicuous placeholder digests, with one declared attempt each. They do not estimate a failure rate on realistic citations, provide a cross-model susceptibility comparison or establish that tool availability caused the observed difference. The tool-assisted condition changes hashing and integer arithmetic together, with separate calls and unrecorded sampling.

\section{Archived procedures and score inventory}
\label{app:procedures}

The archive combines procedures with different computational assistance, context sizes and execution provenance. Table~\ref{tab:supp-procedures} distinguishes their available evidence. Display names identify the supplied runs; no provider API model string, immutable serving snapshot or numeric temperature is available for any of the three foundation-model runs. Configured harness identifiers are not serving attestations. Reproducing a score from saved outputs does not reproduce inference or independently establish the reported restrictions.

\begin{table*}[t]
\caption{Procedure inventory. These are different archived procedures on the same authored key, not a matched comparison of model performance. ``Declared'' conditions and activity self-reports are not independently audited runtime guarantees.}
\label{tab:supp-procedures}
\small
\begin{tabularx}{\textwidth}{@{}p{0.22\textwidth}p{0.40\textwidth}X@{}}
\toprule
Record & Integrity assistance and context & Retained inference evidence\\
\midrule
Bounded checker & Code-based identity, digest, grammar and integer checks; capability ablations share inputs & Checker code, predictions and score reports\\
Claude Fable 5.1, no compute & Declared per-episode fresh instance; prompt reading and response writing; no computational tools & Prompt wrappers, literal responses and activity self-reports\\
Claude Fable 5.1, hash/arithmetic & Declared per-episode fresh instance; Bash restricted by instruction to SHA-256 and integer arithmetic & Prompt wrappers, literal responses and activity self-reports\\
GPT-6.1 Sol Ultra & Frontend supplies citation presence, kind, version and digest comparisons; four primary 17-episode contexts & Instructions, literal responses, integrity observations and worker self-reports\\
Grok 4.7, supplied predictions & Supplier describes per-episode judgment with hashing and exact-replay tools & Structured predictions and procedure account; native prompts and execution traces absent\\
\bottomrule
\end{tabularx}
\end{table*}

\subsection{GPT-6.1 Sol Ultra}
The primary text run covers sixty reporting episodes and eight diagnostics in four declared fresh contexts, each containing seventeen episodes. Their reporting/diagnostic compositions are 17/0, 14/3, 15/2 and 14/3. In total, 37 of the sixty reporting episodes share a context with at least one same-family sibling. The sixteen sibling groups comprise eleven pairs and five triples across all twelve reporting families. No diagnostic siblings co-occur. These counts establish opportunities for within-context comparison, not evidence that the model used them or a measurement of their effect.

The frontend removes unaided digest computation from this procedure. It supplies no source interpretation, policy recommendation or claim label. Applying an observed citation mismatch as a prerequisite for grounding remains a model task. Thus supplied integrity observations do not make the response-A failure logically impossible.

The run was organised using the artifact-authoring assistant system. Archived worker instructions declare no inherited conversation history and prohibit memory or project lookup; parent configuration was inherited. Isolation from persistent memory, project context and instruction state was not independently verified. Shared construction and auditing provenance therefore adds dependence beyond the generator/checker conformance relationship. One attempt per batch is declared, without repair. The failure policy maps missing or invalid batches to U placeholders that are counted incorrect even against reference U, with fixed denominators. No batch failures are reported.

The exact display name GPT-6.1 Sol Ultra was supplied after the run. Historical labels in archived inputs remain unchanged. This naming disclosure does not supply an immutable serving identity or alter any response.

\subsection{Claude Fable 5.1}
The archive contains 136 audit calls: sixty reporting episodes and eight diagnostics under each of two declared computational conditions. Fifty additional calls perform source-only extraction, one per document and modality, giving 186 calls with no recorded response failures. The configured identifier is \texttt{claude-fable-5-1}, with declared harness parameter \texttt{fable}. The supplied run account explicitly notes that a serving-model fallback cannot be verified from inside the session. Sampling used harness defaults; temperature was not recorded.

Each call is described as a fresh instance receiving only its exported prompt. The operator had already seen reference labels and checker code. Tool restrictions and reference isolation were instruction-based; response-level activity reports do not independently attest execution. No-computation calls still read their prompt and write their response, so the condition does not claim absence of all filesystem interaction. The additional Bash condition permits SHA-256 and integer arithmetic. Only this model has an archived condition without computational tools; equivalent conditions for the other two systems are absent.

The original pending metadata remains in the frozen release as historical information. The current Claude evidence is its supplied-results record, extracted response bundle and locally rescored reports. Saved prompt manifests allow comparison of per-call prompt bytes. The retained chronology, attempt count and claimed isolation are declarations, not properties established by archive timestamps.

\subsection{Grok 4.7}
The supplied structured predictions cover the hundred original episodes and eight diagnostics. The reporting and diagnostic portions are rescored separately. The procedure account describes per-episode reasoning with hashing and exact replay, but native prompts and execution traces were not supplied. Consequently the prediction labels can be rescored, while the inference procedure cannot be examined. No claim-level predictions have been reconstructed from aggregate scores.

\subsection{Separate reporting and diagnostic counts}
Table~\ref{tab:supp-counts} inventories conformance counts with fixed denominators. Agreement is the number of labels equal to the authored key. Resolution counts S or C predictions. FS/S is false support divided by all S predictions. False support is an endorsement whose reference is C or U; it is not evidence that the underlying passage or merchant decision is factually false.

\begin{table*}[t]
\caption{Agreement with the same authored key under different procedures. Reporting and diagnostics remain separate. Shared semantics, inspected fixtures and unmatched procedures preclude interpreting this inventory as independent validation or a model ranking. FS/S is false support per endorsement.}
\label{tab:supp-counts}
\small
\begin{tabularx}{\textwidth}{@{}Xrrrrrr@{}}
\toprule
 & \multicolumn{3}{c}{Reporting: 420 claims} & \multicolumn{3}{c}{Diagnostics: 56 claims}\\
\cmidrule(lr){2-4}\cmidrule(l){5-7}
Record & Agreement & Resolved & FS/S & Agreement & Resolved & FS/S\\
\midrule
Checker: execution & 264/420 & 232 & 0/193 & 40/56 & 32 & 0/20\\
Checker: +identity & 322/420 & 290 & 0/247 & 48/56 & 40 & 0/28\\
Checker: +grounding & 418/420 & 386 & 0/331 & 56/56 & 48 & 0/36\\
Claude Fable 5.1, no compute & 359/420 & 331 & 2/277 & 43/56 & 35 & 0/23\\
Claude Fable 5.1, hash/arithmetic & 420/420 & 388 & 0/333 & 56/56 & 48 & 0/36\\
GPT-6.1 Sol Ultra & 420/420 & 388 & 0/333 & 56/56 & 48 & 0/36\\
Grok 4.7, supplied predictions & 420/420 & 388 & 0/333 & 56/56 & 48 & 0/36\\
\bottomrule
\end{tabularx}
\end{table*}

The three assisted foundation-model records reach the reporting key's resolution ceiling of 388/420 and the diagnostic key's ceiling of 48/56. This key cannot distinguish them by agreement. The full checker abstains on two reporting commitments outside its grammar. The no-computation Claude run has 61 reporting disagreements: 54 lineage S-to-U, one lineage C-to-U, four dispute-grounding S-to-U and two commitment-grounding U-to-S. Its thirteen diagnostic disagreements are all S-to-U: eight lineage and five dispute-grounding claims. Capability-limited abstention is distinct from an unjustified endorsement, even though both reduce authored-key agreement.

\section{Paired presentation and source-only extraction}
\label{app:modalities}

Two paired experiments use different model roles. The GPT-6.1 Sol Ultra rendered-evidence audit retains populated records and frontend integrity observations, replacing source passages with PNG cards. Its text control and image condition each use twelve episodes in two declared fresh six-episode batches. Each matches 84/84 authored labels, including 24/24 grounding labels, with 68 S, seven C, nine U and no false support among 68 endorsements. These are paired views of the same audit episodes, not an independent visual dataset.

The Claude Fable 5.1 source-only task instead withholds consumed audit inputs and policy records. Twenty distinct contract texts each require commitment and framework duration; five distinct transaction summaries each require an observation window, eligible count and disputed count. This gives 55 typed fields on 25 unique documents per modality. Reference fields are parsed from authored templates. Each document appears once as text and once as a mildly degraded rendered image. Extracted values do not replace the originally injected inputs of the audit suite.

\begin{table}[t]
\caption{Claude Fable 5.1 source-only extraction on the same 25 documents in two presentations. Quote matching and typed-field correctness are distinct checks.}
\label{tab:supp-extraction}
\small
\begin{tabularx}{\columnwidth}{@{}Xrr@{}}
\toprule
Measure & Text & Image\\
\midrule
All typed fields correct, documents & 25/25 & 25/25\\
Correct typed fields & 55/55 & 55/55\\
Literal quote substring match & 20/25 & 21/25\\
Whitespace-normalised quote match & 25/25 & 24/25\\
Response failures & 0 & 0\\
\bottomrule
\end{tabularx}
\end{table}

Table~\ref{tab:supp-extraction} separates value extraction from quotation fidelity. Literal matching tests whether the emitted supporting quote is an exact character substring of the source text. Whitespace normalisation collapses whitespace before substring comparison. One image response substitutes a separator for source line breaks, so its fields are correct but its quote fails the whitespace-only comparison. No framework-as-commitment extraction occurs in these saved responses. Ceiling counts on these templates do not establish realistic document or integrated-pipeline performance.

The rendering supplement supplies the exact renderer, fixed seed 20261004 and unchanged source-document records used to regenerate the Claude PNGs. Rendering uses a fixed 1240-by-760 canvas, DejaVu fonts, mild rotation, Gaussian blur and speckle noise. Source order is part of the random-number sequence. A recorded verification with Python 3.12.14, Pillow 12.3.0 and FreeType 2.14.3 reproduced all 25 archived PNGs byte for byte and matched every decoded image. This verifies regeneration of the images; it neither makes them realistic scans nor independently verifies historical runtime image delivery. The source-record JSONL includes target fields and is reproduction material, not a blinded inference input.

\section{Data, code and computational record}
\label{app:artifacts}
\label{sec:supp-reproducibility}

The accompanying ancillary material retains frozen release 1.2.2 in \path{anc/explainability-boundaries-v1.2.2.zip.b64}, with separate rendering and run-provenance supplements. Base64 encoding preserves the original ZIP bytes during submission processing. The script \path{anc/restore_release.py} decodes the archive, verifies its SHA-256 digest and extracts the \texttt{explainability-boundaries} directory. File paths below are relative to that directory. The original main-corpus version identifier remains 1.0.0 and the separate diagnostic identifier remains 1.1.0. Corpus and protocol versioning do not imply additional inference. No model is invoked when the archived results are rescored.

Within that release, \path{data/cases.jsonl} contains observable main episodes; \path{data/reference.jsonl} contains authored labels and rationales; and \path{data/metadata.json} records family, split and construction membership. The corresponding diagnostic files are separate under \path{data/diagnostics/}. \path{boundarybench/audit.py} implements the seven checks, while the evaluation and model-adapter modules validate coverage, convert saved outputs and score labels. \path{scripts/generate_data.py} regenerates deterministic fixture files. The accompanying unit tests check boundary conditions, polarity, coverage, integrity and conversion behaviour.

The GPT archive is retained under \path{data/pilot/} and \path{results/pilot/}; the supplied Claude and Grok materials are under \path{results/external/}. Fixed prompts, literal responses where available, converted predictions, raw source images, failure records and score reports retain their separate roles. The scorer does not score rationale quality. A correct final label therefore does not establish that every evidence prerequisite was actually checked in the generating process.

The core checker, generator and offline scorer require Python 3.10 or later and the standard library. No credentials, model service or network access are required for archived-label rescoring. Image regeneration additionally requires the pinned rendering environment. The separate run-provenance script joins existing archive metadata and records input hashes; it does not judge claims. Manifests identify exact file bytes, while the accompanying image-verification record captures runtime and font hashes. Byte reproducibility establishes consistency of a computational record, not independent confirmation of the audit contract or its practical benefit.

All merchants, documents, counts, decisions and histories in the suite are fictional. Neither confidential business records nor personal banking data underlie these results. The suite is public construction material, including its answer key, and is unsuitable as a secret benchmark. Independently authored documents, independent blinded labels, matched model procedures, repeated runs and a reviewer-benefit study remain outside the reported evidence.


\par\smallskip
{\small\noindent\textit{Acknowledgements.} This research has been supported by the FINAI Project Malta Enterprise R\&D Grant.\par}

\end{document}
